% Generated draft: requires mathematical and source review.
\documentclass[11pt]{article}
\usepackage[margin=25mm]{geometry}
\usepackage{fontspec}
\setmainfont{Latin Modern Roman}
\usepackage{amsmath,amssymb,mathtools,unicode-math}
\setmathfont{Latin Modern Math}
\usepackage{enumitem,needspace,xurl,hyperref}
\setlength{\parindent}{0pt}
\setlength{\parskip}{5pt}
\setlength{\emergencystretch}{3em}
\newcommand{\N}{\mathbb{N}}
\newcommand{\Z}{\mathbb{Z}}
\newcommand{\Q}{\mathbb{Q}}
\newcommand{\R}{\mathbb{R}}
\newcommand{\C}{\mathbb{C}}
\newcommand{\F}{\mathbb{F}}
\newcommand{\A}{\mathbb{A}}
\newcommand{\PP}{\mathbb{P}}
\newcommand{\E}{\mathbb{E}}
\newcommand{\eps}{\varepsilon}
\newcommand{\dd}{\,\mathrm{d}}
\begin{document}
\section*{UnsolvedMath ID: 9400110}
\textbf{Source problem code:} catalogue entry
\subsection{Definitions and mathematical statement}
The source names Pillai's conjecture and asks for finiteness of solutions of
\[
 Ax^m-By^n=C,\qquad(m,n)\ne(2,2).
\]
Its domains and nondegeneracy hypotheses are omitted. A common intended weighted formulation fixes positive integers \(A,B\) and a nonzero integer \(C\), and asserts finiteness of
\[
 \{(x,y,m,n)\in\Z_{\geq2}^4:
 Ax^m-By^n=C,\ (m,n)\ne(2,2)\}.
\]
Here \(A,B,C\) are fixed before solutions are counted, while both bases and both exponents vary. This displayed domain is an explicit interpretation requiring source confirmation, not an unmarked completion of ``for any \(A,B,C\)''.

The restrictions matter. Allowing \(C=0\) includes \(A=B=1\), \(m=n=3\), \(x=y\), with infinitely many solutions despite \((m,n)\ne(2,2)\). Allowing exponent one or base one creates further elementary infinite families. Thus the excerpt literally stated is not yet a precise nondegenerate finiteness conjecture.

The summary reports Catalan's case \(A=B=C=1\), with exponents greater than one, as resolved by Mih\u{a}ilescu in 2002. This is a reported special case, not a solution of the general weighted question. Its phrase ``infinitude of the difference'' is not substituted for the source's finiteness assertion.
\subsection{Short English statement}
For fixed nondegenerate coefficients, must the weighted perfect-power equation have only finitely many solutions outside the simultaneous-square case? The source omits essential coefficient, base, and exponent domains; the standard positive-base formulation is marked as an interpretation.
\subsection{Sources}
\begin{itemize}
\item Local UnsolvedMath registry, permanent ID 9400110, source code catalogue entry; no proposer or proposal year supplied; update 2026-08-06; named attribution Pillai.
\item Wikipedia, \emph{List of unsolved problems in mathematics}, Number theory, Diophantine equations bullet 110, wikitext line 139, pinned revision 1366636767: \url{https://en.wikipedia.org/w/index.php?title=List_of_unsolved_problems_in_mathematics&oldid=1366636767#Number_theory}. Registry access date 2026-07-29; MediaWiki extraction. The item is staged pending semantic deduplication and primary-source status review; redistribution review remains required. Wikipedia attribution and CC BY-SA 4.0 handling must be retained for release.
\item The summary names Wikipedia's ``Pillai's conjecture'' entry and Mih\u{a}ilescu's 2002 Catalan resolution, without exact primary references. These literature reports are not independently verified here.
\end{itemize}
\end{document}
